% Pursuing Stacks (cotes 134-2 to 134-8): archive cote 134-7, site page 46 (PDF page 47), typescript page 488.
% typescript. Draft compared: ps6.tex:2268-2332.
% AI-transcribed by Claude Opus 5.5 (claude-opus-5-5, adaptive thinking, effort high), Message Batches API, batch ps3 (msgbatch_01XdFmuVVBAcgc7EFxdtCqsj), 2026-11-12; not yet checked by eye.
\scanpage{134-7/46}
\mspage{488}

unipotent bundles $X$ we are going to work with will be smooth i.e. $M$ pro-%
jective (and we may even get away with free $M$'s, if we need so). The
natural idea here for getting (20) via (21) in full generality, is to
use a projective resolution of $M$ (even a free one), but I'll not try to
work this out now. The main impression which remains is that for the
more relevant cases (involving cohomology \xout{\unsure{inva}} groups of a \uline{smooth}
unipotent bundle\inserted{$X$} at any rate, with quasicoherent coefficients), the
quasicoherent "pro"-linearization $L(X)$ is just as good for computing
cohomology invariants, as the forbidding $\underline{O}_k^{(X)}$ modules we were shrinking
from.

If now we take an $X_*$ instead of just $X$, namely a ss complex in
$U(k)$, and assuming the components $X_n$ to be smooth (to be safe), the
isomorphisms (20') should give rise to isomorphisms (for $F=W(N)$)
\[ (22) \qquad H^n(X_*, F ) \hand{$\simeq$} \operatorname{Ext}^n_{\underline{O}_k}(L(X_*), F) \hand{$\simeq$} \operatorname{Ext}^n_k(L(X_*),N) , \]
where the $\operatorname{Ext}^n$ should be viewed as Hyperext functors (not term by
term), and where \inserted{in the last member} $L(X_*)$ may be interpreted as a chain complex in
Pro $(\mathrm{Ab}_k)$. The first member of (22) is the kind of group occuring as
obstruction group in the Postnikov-type devissage of $X$\hand{$_*$}into linear
structures $W(M(i)_*)$. The chain-pro complex $L(X_*)$ may still look a
\margin{x} little forbidding, ou\over{t}{r} hope though now is that in the "pointed" case
we are really interested in, we may replace $L(X_*)$ by $L$\deleted{\illegible{1}}$_{pt}(X_*)$, which
is just a true honest chain complex in $\mathrm{Ab}_k$ . Now from (16) we get indeed
a canonical map
\[ (23) \qquad \operatorname{Ext}^n_k(L(X_*), N ) \hand{$\longrightarrow$} \operatorname{Ext}^n_k(L_{pt}(X_*), N ) \]
and we hope that this is an isomorphism, under suitable assumptions on
$X_*$ , the most basic one I can think of now being
\[ X_o = e . \]
The map (23) w\over{\unsure{i}}{a}s \inserted{defined as the} transposed of a map of chain complexes in \deleted{$\mathrm{Ab}_k$} \hand{Pro $\mathrm{Ab}_k$}
\[ (24) \qquad L_{pt}(X_*) \hand{$\longrightarrow$} L(X_*) \quad , \]
\xout{by applying} deduced from (16) by applying it componentwise. We recog-%
nize here,\deleted{\illegible{1}} but with a different interpretation (which seems to me
"the correct" one), the second map in the often referred\hand{-}to diagram (w)
of page 454 , or (V\hand{)} in yesterday's reflections (p.478). To say that
it gives rise to isomorphisms (23), for any $n$ and any \over{\unsure{m}}{M}odule $N$, should
be equivalent to saying that (24) is a quasiisomorphism\deleted{s} - but to
make sure \inserted{I \deleted{expected}} \deleted{\illegible{1}}, \inserted{should demand \deleted{\unsure{a}}} \xout{and prove will} \deleted{amount} a little work on foundational
matters on pro-complexes I guess\hand{;} \hand{Also, to see}

\margin{46}
